\documentclass[11pt]{article}
\usepackage[margin=1.15in]{geometry}
\usepackage{amsmath,amssymb,booktabs}
\usepackage{microtype,xcolor,mdframed}
\usepackage[colorlinks=true,linkcolor=blue!50!black,urlcolor=blue!50!black]{hyperref}
\newmdenv[backgroundcolor=blue!4,linecolor=blue!35,linewidth=0.8pt,
  skipabove=9pt,skipbelow=9pt,innerleftmargin=10pt,innerrightmargin=10pt]{keyidea}
\newcommand{\code}[1]{{\small\ttfamily #1}}
\setlength{\emergencystretch}{1.5em}
\title{How Nearly Repeating Parity Forces a Return:\\
\Large A Friendly Guide to Finite Collatz Defects}
\author{Companion to ``Finite Parity Defects Force Collatz Returns''\thanks{Written with Codex (OpenAI). The technical paper and Lean source give the exact statements and proofs. No claim of literature priority is made.}}
\date{September 5, 2026}
\begin{document}
\maketitle

\section{The result in ordinary language}
Use the shortcut Collatz rule: halve an even number; triple an odd
number, add one, and halve it. Record a $1$ at each odd number and a
$0$ at each even number. Now compare this parity diary with a copy
shifted forward by $q$ places.

If the two copies disagree only a few times over a sufficiently long
finite stretch, the orbit must actually return to an earlier number.
The theorem gives both the required stretch length and a short list
of checkpoints where the return must occur.

Let $K$ satisfy $N+1\le2^K$ for the starting number $N$. For instance,
$N=3$ allows $K=2$. If there are at most $D$ disagreements, the
sufficient number of comparisons is
\[
L=(2^{D+1}-1)(q+K).
\]
Those comparisons use $L+q$ parity letters. The possible return
checkpoints are
\[
s_k=(2^k-1)(q+K),\qquad k=0,1,\ldots,D.
\]
\begin{keyidea}
\textbf{What Lean proves.} At one of these checkpoints, the orbit
value equals the value $q$ steps later. For $q>0$, the future repeats
and the orbit is bounded. For $q=2$ and a positive starting number,
the theorem goes further: the orbit reaches one.
\end{keyidea}
The result allows disagreements anywhere in the checked interval.
It does not assume a particular Sturmian slope or a favorable density
of even steps. The cost is that the sufficient length grows
exponentially with the allowed number of disagreements.

This is a new proved criterion in the repository. We have not proved
that every positive starting number satisfies it.

\newpage
\section{Why an exact interval matters}
Two distinct starting numbers sharing their first $r$ parity letters
differ by a nonzero multiple of $2^r$. Their difference is therefore
at least $2^r$. This is the arithmetic fact behind the earlier exact
repetition criterion.

There is also a simple size bound. After $u$ shortcut steps,
\[
T^u(N)+1\le2^u(N+1)\le2^{u+K}.
\]
It is deliberately generous, but holds for every choice of parities.

At a time $s$, compare the two orbit values $T^s(N)$ and
$T^{s+q}(N)$. If their parity diaries agree for the next $s+q+K$
letters, the divisibility lower bound is too large for these two
values to be distinct. They must be equal.

\begin{keyidea}
\textbf{The local choice.} At every checkpoint $s$, either the orbit
returns after $q$ steps, or there is a parity disagreement before
time $2s+q+K$.
\end{keyidea}
Set $A=q+K$. Choose checkpoints $0,A,3A,7A,\ldots$ so that the
search windows fit together without overlapping:
\begin{center}
\begin{tabular}{lll}
\toprule
Checkpoint&Window in which a defect is forced&Next checkpoint\\
\midrule
$0$&$[0,A)$&$A$\\
$A$&$[A,3A)$&$3A$\\
$3A$&$[3A,7A)$&$7A$\\
$7A$&$[7A,15A)$&$15A$\\
\bottomrule
\end{tabular}
\end{center}
If every one of $D+1$ checkpoints failed to return, these disjoint
windows would contain $D+1$ different disagreements. That contradicts
the allowed total of $D$.

This replaces a direct attempt to bound swap costs. Moving one parity
letter across a late adjacent position can have a very large weighted
arithmetic cost. The window argument instead uses the exact stretches
between disagreements.

\newpage
\section{Allowing edits to a repeating template}
Suppose a finite parity diary differs in at most $e$ letters from a
repeating template with period $q$. Each edited letter can affect at
most two comparisons with the shifted copy: one in which it is the
first letter, and one in which it is the second.

Thus at most $e$ edited letters cause at most $2e$ shift disagreements.
The sufficient comparison length becomes
\[
L=(2^{2e+1}-1)(q+K),
\]
and the template is checked for $L+q$ letters. The edit positions can
be arbitrary. Lean checks this extension too.

For a concrete example, start at $3$:
\[
3,5,8,4,2,1,2,1,\ldots.
\]
Its diary begins $11000101\ldots$. Compare it with its shift by two
places. The only disagreements are at times $0,1,3$. With $D=3$,
$q=2$, and $K=2$, the sufficient comparison length is $60$.
The checkpoints are $0,4,12,28$. Already at time $4$, the value $2$
equals the value two steps later.

The same diary differs from the template $010101\ldots$ only at
times $0$ and $3$. This gives $e=2$ edited letters, whose more
conservative template bound is $L=124$. The distinction matters:
edited letters and shift disagreements are different counts.

\section{Shorter windows from a sharper height bound}
The first theorem allowed one bit of height growth per step. We can
improve this without assuming anything new about the orbit. Keeping
the added-one correction, the shortcut map satisfies
\[
2^u\bigl(T^u(N)+1\bigr)\le3^u(N+1).
\]
Five steps therefore need at most three bits of growth, since
$3^5=243<256=2^8$. Checking the remaining zero to four steps gives
\[
T^u(N)<2^{\lceil3u/5\rceil+K}.
\]
Use that smaller allowance to choose the checkpoints:
\[
\sigma_0=0,\qquad
\sigma_{k+1}=\sigma_k+\left\lceil\frac{3(\sigma_k+q)}5\right\rceil+K.
\]
The same counting proof works with the shorter horizon
$L\ge\sigma_{D+1}$. Edited templates use $\sigma_{2e+1}$.

\begin{keyidea}
\textbf{A checked improvement.} Lean proves that these horizons are
never larger than the original ones. Their geometric growth has an
upper bound with factor $8/5=1.6$, improving the old factor of $2$.
The hypotheses about the seed and the allowed defects are unchanged.
\end{keyidea}
\newpage
For lag two, the worked examples become
\begin{center}
\begin{tabular}{rrrrr}
\toprule
Start $N$&Bit bound $K$&Defects $D$&Old horizon&New horizon\\
\midrule
3&2&3&60&36\\
7&3&5&315&128\\
9&4&6&762&250\\
\bottomrule
\end{tabular}
\end{center}
Lean checked the new horizons and finite defect counts. These illustrate
the theorem; they are not new computational Collatz records.

\section{What remains between this and Collatz}
\begin{keyidea}
\textbf{The missing step.} We know what happens if a finite parity
prefix is sufficiently close to periodic. We do not know that every
hypothetical counterexample must have such a prefix.
\end{keyidea}
For a general positive period, the return could belong to a nontrivial
cycle. For period two, Lean proves that the only positive return values
are $1$ and $2$, so the certificate implies convergence to one.
But its universal availability is still unproved. Every orbit reaching
one eventually alternates; assuming that all orbits have the needed
period-two certificate would simply rephrase the convergence problem.

The shorter windows still require an independent restriction on the
number of defects. No such universal restriction has been established.
The expanded audit checks 44 selected declarations using only Lean's
standard foundational axioms; it is not an independent kernel replay.

The technical paper is \code{paper/defects.pdf}, with source
\code{paper/defects.tex}. The formal development is
\code{lean/Collatz/ParityDefects.lean} and
\code{lean/Collatz/DefectHeight.lean}; the research note is
\code{analysis/PARITY\_\allowbreak{}DEFECTS.md}.
\end{document}
